\clearpage
\section{Yukawa Potential}
\begin{colbox}
	The following potential field is given
  \begin{equation}
    \phi(r) = \frac{1}{r}e^{-Kr}; \quad K>0
  \end{equation}
  \begin{enumerate}
    \item Calcualte the gradient for \(r\neq0\)
    \item Calcualte the laplacian for \(r\neq0\)
    \item Use gauss' law to find an expression for the laplacian while includes the origin
    \item Show that \begin{equation}
      \iiint\dd^3x\laplacian{\phi}=0
    \end{equation}
  \end{enumerate}
\end{colbox}
\subsection{gradient}
The gradient in spherical coordinates is defined as
\begin{align*}
    \Vec{\nabla}\phi &= \partial_r\phi\hat{r}+\frac{1}{r}\partial_\theta\phi\hat{\theta}+\frac{1}{r\sin\theta}\partial_\varphi\phi\hat{\varphi}
\end{align*}
therefore
\begin{align*}
    \Vec{\nabla}\phi &= \left(-\frac{1}{r^2}e^{-Kr}-\frac{K}{r}e^{-Kr}\right)\hat{r}\\
    &= -\left(\frac{1}{r^2}+\frac{K}{r}\right)e^{-Kr}\hat{r}
\end{align*}
\subsection{laplacian}
The laplacian in spherical coordinates is defined as
\begin{equation*}
    \nabla^2\phi = \frac{1}{r^2}\partial_r(r^2\partial_r\phi)+\frac{1}{r^2\sin\theta}\partial_\theta(\sin\theta\partial_\theta\phi)+\frac{1}{r^2\sin^2\theta}\partial^2_\varphi\phi
\end{equation*}
therefore
\begin{align*}
    \nabla^2\phi &= \frac{1}{r^2}\partial_r(-\left(1+Kr\right)e^{-Kr})\\
    &= \frac{1}{r^2}(-Ke^{-Kr}+\left(1+Kr\right)Ke^{-Kr})\\
    &= \frac{1}{r^2}(K^2re^{-Kr})\\
    &= \frac{K^2}{r}e^{-Kr}
\end{align*}
\subsection{Gauss' law}
Gauss' law for this subsection is
\begin{equation*}
    \iiint_V\nabla^2\phi d\boldsymbol{x} = \int_{\partial V}\boldsymbol{\nabla}\phi \cdot d\boldsymbol{a}
\end{equation*}
therefore, using the result from subsection 2
\begin{align*}
    \iiint_V\nabla^2\phi d\boldsymbol{x} &= \int_{\partial V}-\left(\frac{1}{r^2}+\frac{K}{r}\right)e^{-Kr}\hat{r} \cdot d\boldsymbol{a}
\end{align*}
We know that \(d\boldsymbol{a}\) points in the radial direction therefore \(\hat{r}\cdot d\boldsymbol{a} = da\) and thus
\begin{align*}
    \iiint_V\nabla^2\phi d\boldsymbol{x} &= \int_{\partial V}-\left(\frac{1}{r^2}+\frac{K}{r}\right)e^{-Kr} da\\
    &= -\left(\frac{1}{R^2}+\frac{K}{R}\right)e^{-KR} \int_{\partial V}da\\
    &= -\left(\frac{1}{R^2}+\frac{K}{R}\right)e^{-KR} \cdot 4\pi R^2\\
    &= -4\pi \left(1+KR\right)e^{-KR}
\end{align*}
Comparing this to a direct calculation
\begin{align*}
    \iiint_V\nabla^2\phi d\boldsymbol{x} &= \iiint_V\frac{K^2}{r}e^{-Kr} d\boldsymbol{x}\\
    &= \int_0^{2\pi}d\varphi\int_0^{\pi}\sin\theta d\theta\int_0^{R}r^2\frac{K^2}{r}e^{-Kr}dr\\
    &= 4\pi K^2\int_0^R re^{-Kr}dr = \left[\begin{aligned}
        U&=r && U'=1\\
        V'&=e^{-Kr} && V=-\frac{1}{K}e^{-Kr}
    \end{aligned}\right]\\
    &= 4\pi K^2\left(\left.-\frac{re^{-Kr}}{K}\right|_0^R+\frac{1}{K}\int_0^R e^{-Kr}dr\right)\\
    &= 4\pi K^2\left(\left.-\frac{re^{-Kr}}{K}\right|_0^R+\frac{1}{K}\left.\left(-\frac{1}{K}e^{-Kr}\right)\right|_0^R\right)\\
    &= 4\pi K^2 \left(-\frac{Re^{-KR}}{K} + \frac{1}{K}\left(\frac{1}{K}-\frac{1}{K}e^{-KR}\right)\right)\\
    &= 4\pi K^2 \left(-\frac{Re^{-KR}}{K} + \frac{1}{K^2}-\frac{1}{K^2}e^{-KR}\right)\\
    &= 4\pi \left(-KRe^{-KR}+1-e^{-KR}\right)\\
    &= 4\pi \left(1-(1+KR)e^{-KR}\right)
\end{align*}
We notice that the direct calculation gives an aswer that is \(4\pi\) off from the result of Gauss' law, which is likely a result of the fact our laplacian definition did not include the origin. Thus we'll add a correction for the origin:
\begin{equation*}
    \nabla^2\phi = \frac{K^2}{r}e^{-Kr}-4\pi\delta^3(\boldsymbol{r})
\end{equation*}
\newpage
\subsection{integral of laplacian}
\begin{align}
  \iiint\dd^3x\laplacian{\phi} &= \iiint\dd^3x\ins{k^2\frac{e^{-Kr}}{r}-4\pi\delta(\vb{r})}
\end{align}
but
\begin{equation}
  \iiint\dd^3xk^2\frac{e^{-Kr}}{r} = 4\pi K^2 \int_{0}^{\infty}\dd r re^{-Kr}
\end{equation}
and since
\begin{equation}
  \int_{0}^{\infty}\dd r re^{-Kr} = \frac{1}{K^2}
\end{equation}
then
\begin{equation}
  4\pi K^2 \int_{0}^{\infty}\dd r re^{-Kr} = 4\pi
\end{equation}
and the second part
\begin{equation}
  \iiint\dd^3x 4\pi\delta(\vb{r}) = 4\pi
\end{equation}
adding the two parts together
\begin{equation}
  \iiint\dd^3xk^2\frac{e^{-Kr}}{r} = 4\pi K^2 \int_{0}^{\infty}\dd r re^{-Kr} = 4\pi-4\pi=0
\end{equation}